\chapter{Реализация}
\label{ch:realizacia}

Методите от глави~\ref{ch:metod} и~\ref{ch:eksperimenti} са превърнати в
работеща и практически полезна софтуерна разработка. Резултатът е едно хранилище с три крайни
продукта, стъпващи върху обща кодова база: (1)~Python пакетът \texttt{drum\_dynamics}
заедно с локална услуга за извод; (2)~приставка (\eng{plugin}) за
цифровата аудио работна станция (\eng{DAW}) \eng{Reaper}, която възстановява
динамиката на нотите в MIDI редактора ѝ (продуктът с най-значима потенциална
практическа приложимост на работата); и (3)~уеб демо в \eng{Hugging Face Spaces}. Продуктовото
име на инструмента е \textbf{„Dynamics Needed''}, докато \texttt{drum\_dynamics} е
името за импорт на пакета — двете са нарочно различни.

\section{Обзор на архитектурата}
\label{sec:arch-overview}

Системата следва принципа „едно ядро, много приложения'' (фиг.~\ref{fig:reaper_plugin}
показва едното от тях). Цялата логика на извода е капсулирана в клас
\texttt{Engine}, който зарежда три от обучените модели (\eng{LightGBM}, трансформър с
\eng{MDN} глава и трансформър с категорийна глава; раздел~\ref{sec:modeli}) и
предоставя единствен метод \texttt{predict}. И приставката за \eng{Reaper}, и уеб
демонстрацията, и командният ред използват въпросният \texttt{Engine} — така
съвпадението между експерименталните резултати и това, което чува крайният
потребител, е гарантирано по конструкция.

Разделянето е от тип клиент–сървър. Тъй като зареждането на моделите и на
\eng{PyTorch} отнема секунди, ,,двигателят'' (\texttt{Engine}) работи като локален процес
(услугата \texttt{drum\_dynamics.\allowbreak serve}), а всеки клиент общува с него през
кратки заявки. Заявката носи само
\emph{структура и време} (кой глас, кога, стил, темпо), а услугата връща нови
стойности за силата (MIDI \eng{velocity}) на всяка нота, които могат да бъдат 
смесени с оригиналните с някакъв коефициент.

\section{Пакетът \texttt{drum\_dynamics} и услугата за извод}
\label{sec:package}

Пакетът е разделен на подпакети по
предназначение: \texttt{core/} (четене на MIDI, каноничното групиране на
гласовете), \texttt{data/} (извличане на входни данни, изграждане на тензори,
разделяне на данните), \texttt{models/} (моделите и изходните глави),
\texttt{eval/} (метрики и оценка), \texttt{viz/} (визуализация и озвучаване),
\texttt{research/} (анализите от Фаза~0) и \texttt{serve/} (услугата за извод).

\paragraph{Услугата за извод.} Подпакетът \texttt{serve} е самостоятелна услуга
върху стандартната библиотека (без външни уеб рамки). Тя изгражда лек HTTP сървър,
слушащ само на локалния интерфейс (\texttt{127.0.0.1}, подразбиращ се порт
\texttt{8765}), с два маршрута: \texttt{GET /health} връща наличните модели,
жанрове и стилове (като проста проверка за наличността на системата);
\texttt{POST /predict} приема редица от ноти и параметри и връща
предсказаните стойности за силата (листинг~\ref{lst:predict}).
Оркестрацията в \texttt{serve.core} е независима от модела: тя изгражда
структурните признаци, извиква избрания модел и накрая за всяка нота смесва (\eng{blend})
предсказаното с оригинала с коефициент между $0$ и $1$ — при $1{,}0$ изходът е
изцяло предсказания, при $0{,}0$ се пазят оригиналните стойности. Записват се само
избраните от потребителя ноти, а редът на изхода се съпоставя точно към входните
позиции.

\renewcommand{\lstlistingname}{Листинг}
\begin{lstlisting}[caption={Заявка и отговор на \texttt{POST /predict} (съкратено).},label={lst:predict}]
// заявка
{ "model": "mdn", "style": "rock", "temperature": 1.0, "blend": 0.8,
  "beat_type": "beat", "bpm": 120.0, "time_signature": "4-4",
  "notes": [ {"index": 0, "pitch": 36, "onset_sec": 0.0,
              "velocity": 100, "selected": true}, ... ] }

// отговор
{ "velocities": { "0": 88, "3": 41, "7": 112, ... } }
\end{lstlisting}

\paragraph{Жизнен цикъл.} Понеже услугата е дълготраен процес, тя се самозакрива
чрез „пазач'' (\eng{watchdog}): излиза при неактивност над зададен праг (по
подразбиране 30~минути) или когато процесът родител (приставката) изчезне — така не
остават процеси ,,сираци''. Теглата на моделите могат да се изтеглят автоматично от
\eng{Hugging Face Hub} при първо стартиране, ако липсват локално.

\section{Приставка за \eng{Reaper} — „Dynamics Needed''}
\label{sec:reaper}

\subsection{Защо чрез ReaScript, а не самостоятелен VST/AU}
\label{sec:why-reascript}

Първоначалната идея за класическа приставка (\eng{VST/AU}) беше изоставена по
концептуална причина: рамките за създаване на приставки като JUCE обработват
аудио и MIDI поток в \emph{реално} време, но не позволяват на приставката да \emph{пренапише}
избраните ноти на място в редактора на партитурата. А без да имат пълният контекст напред във времето,
няма как моделите да изведат стойностите на силата. Освен това първият подход дава възможност на потребителят
и да избере произволен набор от ноти, които да бъдат преоценени.
\eng{Reaper} предоставя пълен скриптов интерфейс към
MIDI съдържанието на такт, поради което избраният носител е
\eng{ReaScript} на \eng{Python}, работещ под вградения в \eng{Reaper} интерпретатор.
Изборът е за конкретен \eng{DAW}, но подходът е представителен: който и да е
\eng{DAW} със сходен скриптов достъп би могъл да бъде обслужен от същия двигател.

\subsection{Архитектура}
\label{sec:reaper-arch}

Приставката е разделена на три слоя така, че по-голямата част от кода да може да бъде
тествана извън \eng{Reaper}:
\begin{itemize}
  \item \textbf{Reaper слой:} Модулът \texttt{dynamics\_needed.py} е
        единственото място, което борави с \eng{Reaper} API и графичната рамка
        (\eng{ReaImGui}). Тук важат ограниченията на вградената среда: само
        \eng{ASCII}, не съществува \texttt{\_\_file\_\_}, извикванията на \eng{Reaper} API
        връщат кортежи, т.н.
  \item \textbf{Ядро (клиент):} Модулите \texttt{dn\_\allowbreak core.py} (изграждане на
        заявката, съпоставяне на жанр към стил, разрешаване на целевите ноти),
        \texttt{engine\_\allowbreak client.py} (стартиране на системата и правене на HTTP заявки),
        \texttt{predict\_\allowbreak worker.py} (фонов процес) и \texttt{boot\allowbreak strap.py}
        (първоначално сваляне) не съдържат \eng{Reaper} API и затова се покриват с
        обикновени \eng{unit} тестове под \texttt{pytest}.
  \item \textbf{Локален сървър:} Услугата от раздел~\ref{sec:package}, стартирана
        като отделен процес и достижима на локалния порт.
\end{itemize}

\subsection{Панел и работен поток}
\label{sec:reaper-panel}

Панелът (фиг.~\ref{fig:reaper_plugin}) е реализиран чрез \eng{ReaImGui} —
графична библиотека за приставки към Reaper, която се
преизчертава на всеки кадър чрез механизма \texttt{defer}.
Контролите отразяват параметрите на модела:
\begin{itemize}
  \item \textbf{Жанр и стил.} Списъкът със стилове се филтрира според избрания
        жанр. Полето „стил'' е активно само за \eng{LightGBM}; трансформърните
        глави (\eng{MDN}, категорийна) се обуславят само по жанр.
  \item \textbf{Модел.} Превключвател между \eng{LGBM}, \eng{MDN} и категорийния.
  \item \textbf{Температура и микс.} Плъзгачът за температура (Temp) е активен само за
        \eng{MDN}; плъзгачът Blend управлява коефициента за смесване на резултатите
        с оригиналните стойности, описан по-подробно в раздел~\ref{sec:package}.
  \item \textbf{„Fill?''} Обозначава дали пасажът е само брейк, тъй като типът на изпълнението
        е вход на моделите.
\end{itemize}
Под контролите се чертае малка графика за бърз преглед на резултантните стойности.
За всяка целева нота тя има два стълба
(текущата в сиво и предсказаната в синьо), така че ефектът се вижда още преди да е
приложен.
Ако няма маркирано подмножество от ноти, по подразбиране се работи с всички ноти
в отворения в MIDI редактора сегмент.

\begin{figure}[H]
  \centering
  \includegraphics[width=\textwidth]{fig_reaper_plugin.png}
  \caption{Приставката „Dynamics Needed'' в \eng{Reaper}. Вляво: прозорецът на приставката.
  Вдясно: MIDI редакторът след натискане на „Apply'' —
  лентата за \eng{velocity} долу вече показва разнообразна динамика (обозначена с различни цветове по нотите)}
  \label{fig:reaper_plugin}
\end{figure}

\subsection{Преглед без блокиране на интерфейса}
\label{sec:reaper-live}

Изводът не бива да „замразява'' интерфейса. Затова заявките се изнасят във фонов режим
(\texttt{predict\_worker.py}) на отделна нишка. На всеки
кадър обвивката образува \emph{подпис на промяната} от текущите параметри и от
структурата на избраните ноти; предсказване се пуска само когато подписът се
промени (при включенo „Live Preview''). Процедурата изчаква кратък интервал
($0{,}15$~s) и изпълнява
само последната заявка, изхвърляйки изпреварените — така бързото местене на
плъзгач поражда едно предсказване, а не поредица. Резултатът се запазва и се
прочита от нишката на интерфейса на следващия кадър. При неуспех се пази последният
успешен резултат и се показва статус, вместо приставката да прекъсне.

\subsection{Спазване на строгото условие вътре в приставката}
\label{sec:reaper-no-leakage}

Приставката прочита от нотите само \emph{структурни} полета (височина, начален
момент, дали са избрани за редакция) и извежда темпото и тактовия размер от проекта — никога
не подава текущата сила \eng{velocity} като вход за предсказването ѝ, в съответствие с
раздел~\ref{sec:no-leakage}. При прилагане („Apply'') всяка засегната нота се
препрочита и се пренаписва \emph{само} полето \eng{velocity}, а всичко останало
(позиция, височина, избор, заглушаване) се запазва непроменено. Действието е
обвито в един блок за отмяна (\eng{undo}), така че потребителят може да го върне с
едно действие. Последно избраните параметри се пазят на ниво пътечка в проекта, за
да се възстановят при следващо отваряне.

\subsection{Жизнен цикъл на сървъра}
\label{sec:reaper-engine}

При отваряне на панела приставката проверява дали сървърът вече е жив чрез
заявка \texttt{GET /health}; ако не е, го стартира и изчаква той да отговори, показвайки
„Starting engine\ldots''. Изчакването е с щедър праг, тъй като първото стартиране
на замразения пакет под \eng{macOS} задейства проверка на сигурността на всички
вложени библиотеки и може да отнеме над минута (следващите стартирания са за
секунди). Стартирането на сървъра протича на фонова нишка, за да не блокира интерфейса.

\section{Инсталация и разпространение}
\label{sec:distribution}

Крайният потребител не бива да инсталира ръчно нито \eng{Python},
нито зависимостите на Python пакета.
Дистрибуцията е по две направления:

\paragraph{Скриптът — през ReaPack.} Самата приставка (няколко \eng{Python}
файла) се публикува през мениджъра на пакети \eng{ReaPack}: потребителят добавя
едно хранилище по URL и инсталира „Dynamics Needed'' от списъка. Индексът
(\texttt{index.xml}) се генерира програмно и изброява главния скрипт и помощните
модули.

\paragraph{Двигателят — замразен, свалян при първо пускане.} Тъй като двигателят
носи \eng{PyTorch}, \eng{LightGBM} и \eng{scikit-learn}, той се пакетира отделно.
С \eng{PyInstaller} се „замразява'' самостоятелен изпълним пакет за всяка ОС (с
процесорна версия на \eng{PyTorch}, за да се събере под ограничението за размер на
издание), към който се добавят теглата на моделите. При първо стартиране
приставката сваля подходящия за платформата пакет от \eng{GitHub Releases},
проверява го по контролна сума \eng{SHA-256}, разархивира го под ресурсната
директория на \eng{Reaper} и (под \eng{macOS}) премахва дублиращите се копия на
\texttt{libomp} чрез символни връзки, за да избегне конфликт при зареждане.
Свързването на версиите (фиксираната версия в приставката, версията в индекса на
\eng{ReaPack} и етикетът на изданието) става едновременно, а старите версии на
двигателя се изчистват. Изграждането на изданията за трите операционни системи е
автоматизирано чрез Github Actions.

\paragraph{Режим за разработчик.} При работа от хранилището се записва
алтернативна конфигурация (интерпретатор от виртуалната среда и коренова папка на
репозиторито), която прескача сваленето и пуска двигателя направо от изходния код —
удобно при разработка и при ръчния тест в самия \eng{Reaper}.

\section{Уеб демо (Hugging Face Spaces)}
\label{sec:demo}

За да е достъпен методът и без инсталация, е публикувана уеб демонстрация в
\eng{Hugging Face Spaces}, изградена с \eng{Gradio}. Потребителят качва MIDI
файл с канал за ударни инструменти (или избира готов пример), задава музикалния контекст (стил,
темпо, тактов размер, тип на такта) и получава обратно MIDI файл с
предсказаната динамика, съпоставено изображение „преди/след'' (матрица от ноти по
време и инструмент, в която цветът кодира нотната сила) и \emph{две} аудио визуализации през звуков
шрифт по стандарта \eng{General MIDI} — оригиналната и обработената — така че
разликата да се сравни както на око, така и по слух. При качване на файл с
разпознаваемо име музикалният контекст се попълва автоматично.

Демонстрацията ползва същия клас \texttt{Engine} като приставката, поради
което поведението съвпада. Работи на хардуера \eng{ZeroGPU} — единствената
безплатна възможност за \eng{Gradio} среда; тъй като моделите са малки и работят добре
на процесор, изводът никога не заявява използване на графична карта. Изходният код живее в
хранилището (папка \texttt{ml/demo/}); скрипт за внедряване сглобява самостоятелен
пакет (приложението, копие на пакета \texttt{drum\_dynamics}, готовите за зареждане
файлове с тегла и звуковия шрифт) и го публикува — теглата и шрифтът не се пазят в
хранилището, а се добавят при внедряване.

\section{Достъпност и възпроизводимост}
\label{sec:availability}

Всички компоненти са публично достъпни. Таблица~\ref{tab:availability} обобщава
входните точки.

\begin{table}[H]
  \centering
  \small
  \begingroup
  \Urlmuskip=0mu plus 3mu\relax
  \renewcommand{\UrlBreaks}{\do\/\do\-\do\.\do\a\do\d\do\y}
  \begin{tabular}{@{}l>{\raggedright\arraybackslash}p{7.8cm}@{}}
    \toprule
    \textbf{Компонент} & \textbf{Достъп} \\
    \midrule
    Изходен код (хранилище) & \url{https://github.com/yalishanda42/dynamics-needed} \\
    Приставка за \eng{Reaper} (ReaPack) & \url{https://yalishanda42.github.io/dynamics-needed/index.xml} \\
    Уеб демо & \url{https://huggingface.co/spaces/yalishanda/dynamics-needed} \\
    Модели & \url{https://huggingface.co/collections/yalishanda/dynamics-needed} \\
    Уебсайт & \url{https://yalishanda42.github.io/dynamics-needed/} \\
    \bottomrule
  \end{tabular}
  \endgroup
  \caption{Публично достъпни компоненти на системата „Dynamics Needed''.}
  \label{tab:availability}
\end{table}
